\chapter{لم تزریق برای زبان‌های مستقل از متن}

همان‌طور که لم تزریق فصل ۸ برای اثبات نامنظم بودن به کار می‌رفت، نسخه‌ای از آن برای زبان‌های مستقل از متن هم وجود دارد. این بار رشته به \textbf{پنج} بخش تقسیم می‌شود و بخش‌های دوم و چهارم \textbf{با هم} تکرار می‌شوند.

\section{صورت لم}

\begin{theorem}[لم تزریق برای زبان‌های مستقل از متن]\label{thm:cfl-pumping}
اگر $A$ زبانی مستقل از متن باشد، عددی مثل $p$ (طول تزریق) وجود دارد به طوری که هر $s\in A$ با $\abs{s}\ge p$ را می‌توان به شکل $s=uvxyz$ نوشت که
\begin{enumerate}
\item برای هر $i\ge0$: \ $uv^ixy^iz\in A$،
\item $\abs{vy}>0$، و
\item $\abs{vxy}\le p$.
\end{enumerate}
\end{theorem}

شرط ۲ می‌گوید دست‌کم یکی از $v$ و $y$ ناتهی است و شرط ۳ می‌گوید $v$، $x$ و $y$ روی هم حداکثر $p$ نماد دارند (اما این بخش می‌تواند در هر جای $s$ باشد، نه فقط ابتدای آن).

\begin{proofidea}
فرض کنید $G$ گرامری برای $A$ باشد و $s\in A$ رشته‌ای خیلی بلند. درخت تجزیه $s$ باید خیلی بلند باشد، پس مسیری طولانی از ریشه به یک برگ دارد. چون تعداد متغیرها متناهی است، طبق اصل لانه کبوتری یک متغیر مثل $R$ روی این مسیر تکرار می‌شود. زیردرخت $R$ بالایی رشته $vxy$ و زیردرخت $R$ پایینی رشته $x$ را تولید می‌کند.

حالا می‌توانیم زیردرخت $R$ پایینی را با زیردرخت $R$ بالایی جایگزین کنیم (و این کار را تکرار کنیم) تا $uv^ixy^iz$ برای $i>1$ ساخته شود، یا زیردرخت بالایی را با پایینی جایگزین کنیم تا $uxz$ ($i=0$) به دست آید. همه این درخت‌ها درخت تجزیه معتبرند:
\begin{ltrfig}
\begin{tikzpicture}[font=\small,scale=0.95]
  \draw[thick,jzMain] (0,4) -- (-3,0) -- (3,0) -- cycle; \node at (0,4.25) {$T$};
  \draw[thick,jzAccent] (0,2.6) -- (-1.8,0) -- (1.8,0) -- cycle; \node[fill=white,inner sep=1pt] at (0,2.6) {$R$};
  \draw[thick,jzAccent] (0,1.3) -- (-0.8,0) -- (0.8,0) -- cycle; \node[fill=white,inner sep=1pt] at (0,1.3) {$R$};
  \node at (-2.4,-0.3) {$u$}; \node at (-1.3,-0.3) {$v$}; \node at (0,-0.3) {$x$}; \node at (1.3,-0.3) {$y$}; \node at (2.4,-0.3) {$z$};
  \begin{scope}[xshift=7cm]
  \draw[thick,jzMain] (0,4) -- (-3,0) -- (3,0) -- cycle; \node at (0,4.25) {$T$};
  \draw[thick,jzAccent] (0,2.6) -- (-1.2,1.0) -- (1.2,1.0) -- cycle; \node[fill=white,inner sep=1pt] at (0,2.6) {$R$};
  \node at (0,1.35) {$x$};
  \node at (-1.9,-0.3) {$u$}; \node at (1.9,-0.3) {$z$};
  \node at (0,-0.75) {$i=0$: \ $uxz$};
  \end{scope}
\end{tikzpicture}
\end{ltrfig}
\end{proofidea}

\begin{proof}
فرض کنید $G=(V,\Sigma,R,S)$ گرامری برای $A$ باشد و $b$ بیشترین تعداد نماد در سمت راست یک قاعده (فرض کنید $b\ge2$). در هر درخت تجزیه هر گره حداکثر $b$ فرزند دارد. پس حداکثر $b$ برگ در فاصله یک گام از ریشه، حداکثر $b^2$ برگ در فاصله دو گام و حداکثر $b^h$ برگ در فاصله $h$ گام هستند. در نتیجه اگر ارتفاع درخت حداکثر $h$ باشد، رشته تولیدشده حداکثر $b^h$ نماد دارد. برعکس، درخت تجزیه رشته‌ای به طول دست‌کم $b^h+1$ ارتفاع دست‌کم $h+1$ دارد.

قرار دهید $p=b^{\abs{V}+1}$. فرض کنید $s\in A$ و $\abs{s}\ge p$. چون $p=b^{\abs{V}+1}>b^{\abs{V}}$، ارتفاع هر درخت تجزیه $s$ دست‌کم $\abs{V}+1$ است. از میان درخت‌های تجزیه $s$، درخت $\tau$ را با \textbf{کمترین تعداد گره} انتخاب کنید. بلندترین مسیر $\tau$ از ریشه به یک برگ طول دست‌کم $\abs{V}+1$ دارد، پس دست‌کم $\abs{V}+2$ گره دارد: یک پایانه (برگ) و دست‌کم $\abs{V}+1$ متغیر. طبق اصل لانه کبوتری، در میان \textbf{پایین‌ترین} $\abs{V}+1$ متغیر این مسیر، یک متغیر $R$ تکرار شده است. دو ظهور آن را $R$ بالایی و $R$ پایینی بنامید و $s=uvxyz$ را این‌طور تعریف کنید: زیردرخت $R$ پایینی $x$ را تولید می‌کند و زیردرخت $R$ بالایی $vxy$ را.

\begin{enumerate}
\item هر دو زیردرخت ریشه $R$ دارند، پس می‌توان یکی را جای دیگری گذاشت. جایگزین کردن زیردرخت پایینی با بالایی (و تکرار) درخت‌هایی برای $uv^ixy^iz$ با $i>1$ می‌سازد، و جایگزین کردن زیردرخت بالایی با پایینی درختی برای $uxz$. پس $uv^ixy^iz\in A$ برای هر $i\ge0$.
\item فرض کنید $v$ و $y$ هر دو تهی باشند. آنگاه با جایگزین کردن زیردرخت $R$ بالایی با زیردرخت $R$ پایینی، درختی به دست می‌آید که همچنان $s$ را تولید می‌کند اما گره‌های کمتری دارد. این با انتخاب $\tau$ به عنوان کوچک‌ترین درخت تناقض دارد. پس $\abs{vy}>0$.
\item $R$ بالایی در میان پایین‌ترین $\abs{V}+1$ متغیر بلندترین مسیر است. پس بلندترین مسیر زیردرخت آن حداکثر $\abs{V}+1$ یال دارد (اگر بیشتر بود، مسیر کل درخت هم بلندتر می‌شد). بنابراین این زیردرخت حداکثر $b^{\abs{V}+1}=p$ برگ دارد، یعنی $\abs{vxy}\le p$.\qedhere
\end{enumerate}
\end{proof}

\begin{warning}
مثل حالت منظم، این لم فقط شرط \emph{لازم} است: هر زبان مستقل از متن شرایط آن را دارد، اما زبانی که شرایط را دارد ممکن است مستقل از متن نباشد.
\end{warning}

\section{مثال‌ها}

روش استفاده مثل فصل ۸ است: فرض می‌کنیم زبان مستقل از متن است، رشته مناسب $s$ را انتخاب می‌کنیم و نشان می‌دهیم \textbf{هر} تقسیم $s=uvxyz$ با شرط‌های ۲ و ۳، یک $i$ دارد که $uv^ixy^iz$ از زبان بیرون می‌افتد.

\begin{example}[$a^nb^nc^n$]\label{ex:anbncn}
$B=\{a^nb^nc^n\st n\ge0\}$ مستقل از متن نیست.

فرض کنید $B$ مستقل از متن باشد و $p$ طول تزریق آن. رشته $s=a^pb^pc^p\in B$ را انتخاب کنید. هر تقسیم $s=uvxyz$ با $\abs{vy}>0$ را در نظر بگیرید:
\begin{enumerate}
\item \textbf{هر کدام از $v$ و $y$ فقط یک نوع نماد دارد.} در این صورت $v$ و $y$ روی هم حداکثر دو نوع از سه نماد $a$، $b$ و $c$ را دارند. پس در $uv^2xy^2z$ تعداد دست‌کم یک نماد زیاد شده و تعداد دست‌کم یک نماد ثابت مانده است. پس تعداد $a$ها، $b$ها و $c$ها برابر نیست و $uv^2xy^2z\notin B$.
\item \textbf{$v$ یا $y$ بیش از یک نوع نماد دارد.} در این صورت در $uv^2xy^2z$ نمادها به ترتیب درست نیستند (مثلاً اگر $v=ab$، رشته شامل $abab$ است). هیچ رشته‌ای در $B$ زیررشته‌ای مثل $ba$، $ca$ یا $cb$ ندارد، پس $uv^2xy^2z\notin B$.
\end{enumerate}
هر دو حالت به تناقض می‌رسند، پس $B$ مستقل از متن نیست. (اینجا حتی به شرط ۳ هم نیازی نشد.)
\end{example}

\begin{example}[رشته تکراری]
$D=\{ww\st w\in\{0,1\}^*\}$ مستقل از متن نیست.

فرض کنید $p$ طول تزریق باشد و $s=0^p1^p0^p1^p\in D$. (دقت کنید $0^p1\,0^p1$ که در فصل ۸ به کار رفت اینجا جواب نمی‌دهد، چون با $v=0$ و $y=0$ در دو طرف $1$ اول قابل تزریق است.) هر تقسیم $s=uvxyz$ با شرط‌های ۲ و ۳ را در نظر بگیرید. چون $\abs{vxy}\le p$، بخش $vxy$ حداکثر دو بلوک مجاور از چهار بلوک $0^p$، $1^p$، $0^p$، $1^p$ را می‌پوشاند.
\begin{itemize}
\item اگر $vxy$ کاملاً در نیمه اول $s$ باشد، با $i=2$ طول رشته کمی زیاد می‌شود و نیمه دوم رشته جدید با $1$ شروع می‌شود (چون بخشی از بلوک‌های اول به نیمه دوم منتقل شده)، در حالی که نیمه اولش با $0$ شروع می‌شود. پس رشته به شکل $ww$ نیست. حالت کاملاً در نیمه دوم هم مشابه است: نیمه اول رشته جدید به $0$ ختم می‌شود و نیمه دوم به $1$.
\item اگر $vxy$ از وسط $s$ رد شود، $vxy$ در بلوک‌های $1^p0^p$ میانی است. با $i=0$ رشته $0^p1^j0^k1^p$ به دست می‌آید که دست‌کم یکی از $j$ و $k$ کمتر از $p$ است. چنین رشته‌ای به شکل $ww$ نیست، چون اگر بود هر نیمه باید $0^p$ و $1^p$ کامل داشته باشد.
\end{itemize}
پس $D$ مستقل از متن نیست. (اگر $uv^ixy^iz$ طول فرد داشته باشد، بدیهی است که در $D$ نیست؛ تحلیل بالا برای طول زوج است.)
\end{example}

\begin{note}
این دو مثال نشان می‌دهند ماشین پشته‌ای چه کاری را \emph{نمی‌تواند} انجام دهد: پشته فقط یک «شمارش» یا یک «مقایسه آینه‌ای» را می‌تواند پیگیری کند ($a^nb^n$ و $ww^R$ مستقل از متن‌اند)، اما نه دو شمارش هم‌زمان ($a^nb^nc^n$) یا مقایسه در یک جهت ($ww$).
\end{note}

\begin{summary}
\begin{itemize}
\item لم تزریق برای \lr{CFL}ها: $s=uvxyz$ با $uv^ixy^iz\in A$، $\abs{vy}>0$ و $\abs{vxy}\le p$.
\item اثبات: درخت تجزیه بلند، تکرار یک متغیر روی بلندترین مسیر، جایگزینی زیردرخت‌ها. طول تزریق $p=b^{\abs{V}+1}$.
\item $\{a^nb^nc^n\}$ و $\{ww\}$ مستقل از متن نیستند.
\end{itemize}
\end{summary}
